\section{Plane maps, forbidden minors, and five colours}\label{sec:plane}
We give the combinatorial version of plane embedding. Each edge has two
distinct darts paired by an involution $\alpha$, even for a loop. A rotation
$\sigma$ is a permutation of the darts whose cycles are exactly the cyclic
orders of outgoing darts at the nonisolated vertices. From a dart, crossing
and taking the next dart in that cyclic order applies $\phi=\sigma\alpha$.
Its cycles are called facial walks. A component consisting of one isolated
vertex contributes one empty facial walk. For disconnected graphs we use a
separate sphere for each component until an outer face is chosen.

All three permutations are tables on $[2m]$. Count cycles by taking the least
unmarked member, following at most $2m$ steps and marking its orbit. A
specified embedding thus has a numerical interpretation entirely inside iPRA.
The connection with plane pictures is that a vertex is a small disc and
an edge an untwisted band. The boundary of these discs and bands follows
$\phi$. All embedding claims below refer to this finite representation.

\begin{lemma}[Face tracing and the spherical certificate]\leavevmode\label{lem:sphere}
For a connected nonempty graph with any rotation there is a natural $g$ with
$n-m+f=2-2g$, where $f$ counts its facial walks. For any fixed spanning tree,
restore the other edges one at a time in their input cyclic positions.
Then $g$ is exactly the number of restorations joining two different facial
walks. In particular $g=0$ has a certificate built using only face splits.
\end{lemma}
\begin{proof}
Initially restrict each cyclic order to the tree. A one-vertex tree has one
empty boundary. Adding a leaf and its edge inserts a trip out and back into
the boundary at the parent's indicated corner. It preserves the number of
boundaries, also when the parent was the singleton. Thus the tree has one
face and $n-(n-1)+1=2$.

To insert an edge, mark its two attachment slots in the current boundary
lists. A slot is the corner between two successive local darts, with a
temporary marker for a degree-zero singleton. When both new darts enter
the same slot, their relative order is that in the final rotation, so one
of the intervening strings can be empty. Cut the boundary at the two marks,
and cross-connect the resulting ends via the two opposite darts of the
new edge. On one old boundary, write its cyclic word as $AB$ cut at those
marks. The new cycles consist of $A$ and one new dart, and $B$ and its
mate. On two old boundaries, now words $A,B$, the new cycle contains
$A$, one new dart, $B$, and the other dart. These splices follow directly
by applying $\sigma\alpha$ at the two changed corners. Every other transition
is unchanged. The singleton's empty word gives two one-dart faces on
inserting its first loop. Consequently $m$ increases by one and $f$ increases
by one or decreases by one. In the latter case the displayed characteristic
drops by two. Induction on the number of restorations proves the formula.
\end{proof}
Call a connected rotation \emph{spherical} if it has such a tree-and-splits
construction. This is the combinatorial spherical map convention. The lemma
proves that the choice of tree and insertion order is immaterial, and gives
Euler's formula $n-m+f=2$ as a consequence of the construction. The face
words also describe the disc picture: a thickened tree is a disc, and at
each split a chord is inserted in a complementary disc between its marked
boundary slots, cutting that disc in two. A graph is called planar here if
each nonempty component admits a spherical rotation. All rotations can be
tested by a scan below $\max(2,2m)^{2m}$. If all components share a common
exterior face, the merged face count $F=1+\sum_i(f_i-1)$ gives
$n-m+F=1+c$ for $c$ components, including the empty graph. These formulas
make no quantification over real curves.

\begin{proposition}[Minor closure and dual trees]\leavevmode\label{prop:dual}
Deletion of edges or vertices and contraction of nonloop edges preserve
planarity in this convention. For a nonempty connected spherical map its
dual has one vertex per face, one edge per primal edge, one face per primal
vertex, and is spherical. The dual of the dual recovers the original map.
The complements of spanning trees are dual spanning trees.
\end{proposition}
\begin{proof}
For a nonbridge $e$ select a tree avoiding it and restore $e$ last. Zero
merges for the full rotation imply zero for the rotation just before that
last split. This gives deletion, including a loop. If $e$ is a bridge, its
removal gives two connected components with the inherited rotations.
Joining two such boundary lists by the bridge joins one face on each side
into one (the two-boundary splice above). Hence the defect
$2-(n-m+f)$ of the joined graph is the sum of the component defects. Both
are nonnegative, so if the sum is zero both are zero. This calculation
also covers isolated components via the empty-face convention.

For contraction of $uv$, $u\ne v$, cut the cyclic lists after its two darts,
remove them and concatenate the remaining lists into the rotation about the
merged vertex. Tracing a face then just skips the deleted darts: arriving
where a deleted dart used to leave now continues at the former other end.
A face consisting solely of deleted darts can occur only when $uv$ is a whole
one-edge component, whose contracted vertex has its empty face by convention.
Thus $f$ is unchanged, while $n,m$ each drop by one. Delete loops and surplus
parallel edges afterward when a simple minor is required. To delete a vertex,
delete all its incident edges and discard its isolated component.

For a map with edges, the dual uses the same paired darts and vertex
rotation $\sigma^*=\phi$. Its face permutation is $\sigma^*\alpha=\sigma$.
The original graph is connected precisely when the darts form one orbit
under $\alpha,\sigma$: moves around a vertex and across an edge lift and
project walks. Since $\phi=\sigma\alpha$, the dual has the same orbit.
The Euler characteristic is unchanged. The assertion for a singleton is
direct. Fix a primal spanning tree $T$. Delete edges outside $T$ in any
order. Each deletion is of a nonbridge and, by restoring it last, reverses
a split, merging two distinct faces. On the dual side the names of the two faces incident with the deleted edge are identified. Thus these are successive
contractions of edges with distinct endpoints. After all deletions the
primal tree has one face, so the selected dual edges connect all original
dual vertices. There are $m-(n-1)=f-1$ selected edges, hence they form a
tree. Apply the same construction to the dual for the converse.
\end{proof}

\begin{theorem}[Euler bounds and explicit obstructions]\leavevmode\label{thm:densplane}
In a connected simple spherical map on $n\ge3$ vertices, $m\le3n-6$.
If it is triangle-free, $m\le 2n-4$. In particular $K_5$ and $K_{3,3}$ are
nonplanar, as is any graph having one of them as a minor or a subdivision
subgraph.
\end{theorem}
\begin{proof}
Each dart is on exactly one facial walk, so face lengths sum to $2m$.
A face of length one needs a loop. A face of length two either uses parallel
edges or immediately reverses a single edge at both ends. Immediate reversal
in the face trace means that the endpoint has degree one. In the latter case
the entire component would be $K_2$. A length-three closed walk in a simple
graph uses three different vertices and forms a triangle. Therefore
$2m\ge3f$ or, under triangle-freeness, $2m\ge4f$. Substitute $f=2-n+m$.
Now $K_5$ has ten edges against a bound of nine. The bipartite $K_{3,3}$
has nine against a triangle-free bound of eight. Apply minor closure.
In a subdivision subgraph, contracting along each subdivided edge recovers
the original graph, giving the last assertion.
\end{proof}

\begin{theorem}[Constructive five-colour theorem]\leavevmode\label{thm:five}
A finite simple graph with spherical rotation certificates has a computable
proper five-colouring.
\end{theorem}
\begin{proof}
Every nonempty such graph has a vertex of degree at most five. In a component
with at least three vertices a contrary degree sum would give $2m\ge6n$,
contrary to Theorem~\ref{thm:densplane}. Components of size one and two
are immediate. For a vertex $v$ of degree at most four, delete it, colour
recursively and use a missing colour. At degree five consider its five
neighbours. If every two were adjacent, their induced $K_5$ would be a
nonplanar subgraph. Search all ten pairs and select two nonadjacent
neighbours $p,q$. Contract the two edges of $p v q$ to a single vertex
$w$ and simplify. This is a smaller planar graph. From its colouring,
transfer the colour of $w$ to both $p,q$ and copy the colours of all
vertices outside $\{p,v,q\}$. Every edge away from $v$ is proper: the
only possible edge lost between two differently named but equally coloured
endpoints would be $pq$, which was tested absent. Among the five neighbours
of $v$ at most four colours now occur. Give $v$ the least palette colour missing
there. All recursion decreases the vertex count. Store the contraction
and relabelling records and unwind at most $n$ frames. For induction fix
$n$ and conjoin correctness for smaller vertex codes and their rotations.
A simplified intermediate graph has at most $\binom n2$ edges.
\end{proof}

\begin{corollary}[Four colours for triangle-free spherical graphs]\leavevmode\label{idx:plane-four-colours-for-triangle-free-spherical-graphs}
Every simple triangle-free graph with spherical rotation certificates is
constructively four-colourable.
\end{corollary}
\begin{proof}
In each component on $n\ge3$ vertices, $2m\le 4n-8$ gives a vertex of
degree at most three. Smaller nonempty components have such a vertex too.
Deleting that vertex preserves both hypotheses. Record $n$ such removals,
then colour the vertices in reverse order. On inserting a vertex at most
three colours are forbidden by previously inserted neighbours.
\end{proof}
